\section{Tổng quan nghiên cứu}
Định hướng tiếp cận truyền thống trong việc kiểm soát chất lượng không khí trong nhà (IAQ) phần lớn phụ thuộc vào các hệ thống thông gió cơ khí, thiết bị màng lọc tĩnh điện hoặc hệ thống điều hòa không khí. Mặc dù các giải pháp kỹ thuật này là nền tảng tất yếu, chúng thường đi kèm với chi phí năng lượng đáng kể, yêu cầu bảo trì liên tục và chưa giải quyết được nhu cầu thiết lập một hệ sinh thái không gian sống xanh trong bối cảnh kiến trúc hiện đại. Dưới lăng kính của công nghệ sinh thái, các phân tích chuyên sâu về hệ thống thủy canh tích hợp Internet vạn vật (IoT) đã minh chứng rằng, thực vật có thể được số hóa để trở thành một cấu phần chủ động trong mạng lưới giám sát vi khí hậu, nơi các biến số như nồng độ CO2, nhiệt - ẩm độ, quang năng và đặc tính dinh dưỡng được lượng hóa theo thời gian thực \cite{hydroponic_iaq_survey}. Xét riêng về mặt thực vật học ứng dụng, lan ý (\textit{Spathiphyllum}) được định vị là loài sinh vật chỉ thị thích hợp nhờ biên độ chịu đựng tốt với ánh sáng khuếch tán, yêu cầu môi trường ấm ẩm và khả năng sinh trưởng bền vững trong môi trường thủy canh \cite{peace_lily_care_bhg,peace_lily_hydro_brightlane}.

Trong kỷ nguyên của nhà thông minh (smart home), các hệ thống tự động hóa hiện hành chủ yếu vận hành xoay quanh trục điều khiển thiết bị chấp hành (chiếu sáng, thông gió, an ninh). Một bộ phận các hệ thống tiên tiến có tích hợp module cảm biến nhiệt - ẩm và CO2; tuy nhiên, luồng dữ liệu môi trường này phần lớn được xử lý như một tham số độc lập, thiếu đi sự tham chiếu và liên kết với bài toán sinh thái học của thực vật. Ở chiều ngược lại, các hệ thống nông nghiệp chính xác (precision agriculture) lại chỉ tập trung vào sinh lý cây trồng (độ ẩm giá thể, pH, độ dẫn điện TDS, quang chu kỳ) mà bỏ qua sự tương tác mật thiết của chúng với chất lượng không khí xung quanh. Sự phân mảnh trong tư duy thiết kế này đã làm suy giảm tiềm năng của mô hình thực vật hỗ trợ IAQ, bởi lẽ động lực học vùng rễ, trạng thái cân bằng dinh dưỡng và sức khỏe tổng thể của cây có mối tương quan tuyến tính với khả năng trao đổi khí và cải thiện vi khí hậu của chúng.

Trên bình diện công nghệ, kiến trúc IoT cung cấp bộ khung hạ tầng cho phép hội tụ các module cảm biến, thiết bị chấp hành, vi điều khiển cục bộ và nền tảng điện toán đám mây thành một cơ chế giám sát thời gian thực liền mạch \cite{wang_smart_agriculture,liu_realtime_iaq,singh_iot_air_quality}. Khi được tinh chỉnh cho mô hình lan ý thủy canh, hạ tầng IoT có năng lực lượng hóa sự sai lệch của các tham số môi trường và dinh dưỡng (ánh sáng, nồng độ ion, pH, nhiệt - ẩm độ), qua đó thiết lập các tập luật cảnh báo và kích hoạt chu trình điều khiển vòng kín. Theo các nghiên cứu thực nghiệm sinh lý học trên \textit{Spathiphyllum}, điều kiện tối ưu trong giai đoạn thích nghi bao gồm độ dẫn điện (EC) dung dịch ở mức 1.2 dS/m, pH xấp xỉ 5.8, mật độ quang thông (PPF) trong khoảng 70--100 $\mu$mol m$^{-2}$ s$^{-1}$ và giới hạn nhiệt độ quanh mức 25~$^\circ$C \cite{peace_lily_ec_ppf}. Song song đó, các cẩm nang hướng dẫn ứng dụng thủy canh vi mô khuyến nghị duy trì dải pH từ 6.0--6.5 và kiểm soát nồng độ TDS định kỳ \cite{peace_lily_hydro_brightlane}. Mặc dù lý thuyết đã được khẳng định, các nguyên mẫu hệ thống hiện hành vẫn chưa giải quyết triệt để bài toán đồng bộ hóa luồng điều khiển hai chiều giữa giao diện người dùng cuối, lõi xử lý IoT và cụm thiết bị viền (edge devices).

Xét về hạ tầng truyền thông, giao thức Wi-Fi thường là ưu tiên mặc định trong không gian dân dụng nhờ tính sẵn sàng cao. Dẫu vậy, việc bố trí các node cảm biến tại các khu vực rìa (ban công, cửa sổ) hoặc tại các không gian bị phân tán bởi kết cấu vật lý thường dẫn đến suy hao tín hiệu và mất ổn định liên kết. Ngược lại, kỹ thuật điều chế LoRa thể hiện tính ưu việt vượt trội đối với việc truyền tải các gói dữ liệu telemetry chu kỳ, kích thước nhỏ và không đòi hỏi băng thông rộng. Do đó, kiến trúc mạng của đề tài tận dụng liên kết LoRa ở tầng thấp (node - gateway), trong khi giao phó nhiệm vụ định tuyến dữ liệu lên nền tảng đám mây ThingsBoard cho gateway ESP32.

Trong thiết kế phần mềm, ThingsBoard được khai thác như một nền tảng lõi đảm nhiệm chức năng quản lý danh tính thiết bị, lưu trữ chuỗi thời gian telemetry, cung cấp công cụ phân tích trực quan và kiến tạo cơ chế RPC (Remote Procedure Call) cũng như attribute cho luồng điều khiển song hướng. Ở phía người dùng cuối, một ứng dụng di động phát triển trên nền tảng Flutter sẽ giao tiếp trực tiếp qua các điểm cuối API của ThingsBoard để kết xuất dữ liệu, cấu hình ngưỡng hoạt động và khởi tạo lệnh điều khiển. Cấu trúc thiết kế này đảm bảo nguyên tắc phân rã trách nhiệm (separation of concerns): bóc tách rõ ràng giữa lớp biên (thu thập dữ liệu), lớp mạng (truyền thông không dây), lớp lõi (xử lý dữ liệu đám mây) và lớp ứng dụng (giao diện tương tác).

Dựa trên những cơ sở luận lý và thực tiễn vừa trình bày, nghiên cứu này đặt trọng tâm vào việc thiết kế và hiện thực hóa một nguyên mẫu IoT tích hợp. Nguyên mẫu này không chỉ đảm bảo tính khả thi về chi phí và quy mô triển khai trong căn hộ chung cư, mà còn hiện thực hóa sự đồng bộ giữa hệ thống quản trị vi khí hậu IAQ và mô hình chăm sóc thực vật tự động hóa. Đề tài không dừng lại ở mức độ thu thập và biểu diễn số liệu, mà còn cung cấp khả năng cảnh báo sớm, can thiệp thủ công từ xa và thiết lập nền tảng tham số vững chắc để hướng tới việc triển khai các giải thuật điều khiển thích nghi trong các pha phát triển kế tiếp.